\section{Opportunity Identification (Part 2)}
% this is some introduction about the topic
\subsection{Pre-class Activities}
\subsubsection{Activity 1:  Generating List of Innovation Opportunities}
\textbf{Objective:} \\
The objective of this activity is to guide you in creating opportunities for innovation based on the innovation charter you developed in Week \ref{section:opportunity identification part 1}. After completing this activity, it is expected that you can explore new ideas that align with the objectives and boundaries outlined in the charter.

\vspace{10pt} 

\noindent To create opportunities for innovation based on the innovation charter you drafted in Week \ref{section:opportunity identification part 1}, follow these steps:
\begin{enumerate}
    \item Refer to the innovation charter you developed in Week \ref{section:opportunity identification part 1} as a guide for generating new opportunities.
    \item Generate a minimum of three opportunities for innovation that align with the objectives and boundaries outlined in your innovation charter.
    \begin{enumerate}
        \item Try to think outside the box (without the help of any LLM) and explore multitude possibilities.
        \item Consider various aspects including market needs, technological advancements, customer preferences, or process improvements.
        \item Most important, ensure that the opportunities you come out should within the confine of the broad scope of the innovation charter and have the potential to bring value and impact to your team project.
    \end{enumerate}
    \item Note-down each opportunity.

    While writing down each opportunity, consider the following optional steps, which are highly recommended to facilitate ease of revision.
    \begin{enumerate}
        \item Include a descriptive explanation of the idea and its relationship to the innovation charter.
        \item Consider the potential benefits, challenges, or risks, and identify any necessary resources or support required to pursue each opportunity.
    \end{enumerate}
\end{enumerate}

\subsection{In-class Activities}
\subsubsection{Activity 1: Compiling List of Innovation Opportunities}
During this stage, it is expected that each group member has individually generated several opportunities for innovation based on the individually drafted innovation charter.
\begin{enumerate}
    \item Refer to the innovation charter that the team developed in Week \ref{section:opportunity identification part 1} as a compass for this activity.
    \item Review and evaluate the identified opportunities.
    \begin{enumerate}
        \item Keep in mind that the objective of this screening process is not to select a single "best" opportunity but to eliminate opportunities with low potential for value creation, while focusing opportunities that warrant deeper exploration.
        \item Remember that the choices made during this screening step are not final and can be revisited and revised later in the process as the project progresses and more insights are gained.
        \item A stickers will be provided to help the voting, but you may consider alternative voting strategies if desired.
    \end{enumerate}
    \item Write down the finalized opportunities, along with their prioritization and rationale.
\end{enumerate}